\subsection{Cycle candidates, retained scales, and outer forcing}

Write
\[
 F_\pm(X)=1\pm\bar aX+\bar BX^2,\qquad
 \mathcal P=\{\bar c=\bar m=0,\ \bar B>0,\
           |\bar a|\le2\sqrt{\bar B}\}
\]
on the normalized max-norm face. At $\mathcal P$ one has
$\bar B\ge1/4$. Indeed either $\bar B=1$ or $|\bar a|=1$,
and the latter inequality gives $\bar B\ge1/4$.
The compactified interval $[0,\infty]$ is always given its ordinary
topology through $p\mapsto p/(1+p)$.
We first cover closing orbits, so that an outer orbit can supply a
coefficient inequality. We will extend the cover to the intervening
comparison arcs after proving that all of those arcs return.

\begin{lemma}[Refined candidate space]\label{rw:fibre:candidate-space}
The limiting data of parameter-dominant counted cycles have a compact
refinement with the following properties.
Away from $\mathcal P$ its coordinates are the normalized
coefficients and the two compactified actual endpoints.
Near $\mathcal P$ it additionally retains the secondary direction
$(C,M)$, the secondary scale and actual first-exit tags.
Every point has one of the local descriptions of Sections~7--10,
including their boundary faces. The same compactification can be
applied simultaneously to an outer cycle and any returning upper
arc between fixed positive initial endpoints and its endpoints.
The assertion about such additional arcs is compactness only;
their existence and local classification will be proved below.
\end{lemma}
\begin{proof}
Let $\chi>0$ be the fixed strong-entry constant in
Lemma~\ref{rw:finite:action-order}. For parameter-dominant
cycles, $p\ge\chi/2$, $q\ge\chi$. Take the intersection, as
$r\downarrow0$, of the closures of their data with $0<\rho<r$
in
\[
 \mathcal S\times[\chi/2,\infty]\times[\chi,\infty],
 \qquad
 \mathcal S=\{(\bar B,\bar m,\bar a,\bar c):
          \max(|\bar B|,|\bar m|,|\bar a|,|\bar c|)=1\}.
\]
This is a compact set. Its points satisfy the nonnegative lower
conditions of Lemma~\ref{rw:finite:lower-cones}. The parameter
$\nu$ tends to zero by Proposition~\ref{rw:geometry:flatness}.

Only the corner at $\mathcal P$ needs further coordinates. Put
\[
 \bar d=\bar c+\nu,\qquad
 \lambda=\max(|\bar d|,\sqrt{|\bar m|}),\qquad
 C=\bar d/\lambda,\quad M=\bar m/\lambda^2
 \quad(\lambda>0).
\]
The direction belongs to the compact face
$\mathcal D=\{\max(|C|,\sqrt{|M|})=1\}$. Use the fixed
sections of Lemma~\ref{rw:secondary:routing}. Retain a finite
flag specifying a wedge exit, a secondary exit, or their overlap.
For a wedge intersection retain $H\in[0,H_0]$ and, if
$H+\lambda>0$, the ratio
$\chi_{\rm exit}=H/(H+\lambda)$.
For a secondary intersection retain its actual radial coordinates
$R_{\rm in},R_{\rm out}\in[0,R_{\max}]$ on the sections
$\phi=\pm\lambda L$, together with
\[
 t=1/P,\qquad \widetilde t=1/Q,\qquad
 0\le t,\widetilde t\le1/p_* .
\]
Here $P,Q$ are the auxiliary linear-flow labels of
Lemma~\ref{rw:secondary:tags}, evaluated at those actual section
points; in particular $t=R/\bar x(T)$ on the incoming section.
They are not defined to be $\lambda p,\lambda q$.
Use a fixed placeholder for an inactive coordinate. Allow all the
indicated flags on an overlap, and take the closure of the resulting
actual-data graph in the finite union of these compact products,
also retaining $\lambda\in[0,\lambda_*]$.
Taking the closure retains radial limits without inverting a map
at zero radius. An orbit in an overlap may have two sets of tags;
keeping both copies preserves compactness and suffices for a finite
subcover.

If $H/\lambda$ tends
to infinity, routing is in the wedge and
$|\bar d|/H+|\bar m|/H^2$ tends to zero. For a finite ratio
the first-exit flag is retained: a first wedge exit has
$H/\lambda\ge L$ and hence the fixed wedge inequalities,
whereas a first secondary exit is in the secondary rectangle.
The equality face has the open overlap proved below. Radial escape is
$R=0$, equivalently $1/P=0$ or $1/Q=0$, and is retained in
the secondary infinity charts. The cases of a finite upper
threshold and its repeated limit are retained with their actual
deficits. Proposition~\ref{rw:secondary:exhaustion} classifies the
entire normalized face, including its repeated and weak boundaries,
and excludes another zero corner. For $\lambda=0$ an actual field
uses only the wedge; secondary directions there record limits with
$\lambda>0$. This classification applies to returning upper arcs.

For completeness the primary alternatives follow directly from
the upper linear flow. A finite nonregular pair is $(p_c,q_c)$
with $\bar m>0$. One finite and one infinite endpoint is a
strict real threshold, a repeated threshold with its whole interval
$p\ge P$, or the weak boundary $\bar m=0$; reflection gives
the other orientation. At two infinite endpoints the strict barrier
of Lemma~\ref{rw:finite:strict-barrier} gives
$\bar c^2+4\bar m\le0$. Strict focus, repeated upper root and
the cubic horizontal boundary are covered in Section~8; the joint
repeated center is Section~9. If $\bar m=0$, then $\bar c=0$
and the lower cones put the point in $\mathcal P$.
At a finite lower loss the first root, its double limit, or the
simultaneous pair gives the alternatives of Section~7.
The two-root and one-infinite statements include roots that escape
to infinity, so these faces are not omitted.

For an intermediate arc, take the product with a second set of actual
tags, sharing the coefficients but retaining its own section points
and first-exit flag.
Use independent initial bounds $p\ge p_0>0$, $q\ge q_0>0$ for
this copy; $q_0$ is uniform on the compact regular entrance.
Impose $p\le p^{\rm out}$, $q\le q^{\rm out}$ before
closing the actual-data graph. This remains a closed subset of a
compact product. The independent tags are needed because an
intermediate arc may take a wedge exit although the outer cycle takes
a secondary exit.
\end{proof}

\begin{lemma}[Uniform selection over compact fibres]
\label{rw:fibre:compact-selection}
In a compact metric product $X$ with continuous outer projection
$\pi:X\to B$ to a compact metric space, let
$X_r$ be the actual data with $0<\rho<r$ and put
$X_0=\bigcap_{r>0}\overline{X_r}$. If a finite union
$U$ of prepared open pieces contains the limiting fibre
$X_0\cap\pi^{-1}(\{b\})$, then
an outer neighborhood and a positive original scale cutoff can be
chosen so that every actual datum over that neighborhood belongs to
$U$. If this construction applies at every outer limit, finitely many
such choices give common cutoffs and bounds.
\end{lemma}
\begin{proof}
Choose $V$ with $\overline V$ disjoint from the compact set
$\pi(X_0\setminus U)$. If no scale cutoff worked, excluded actual data
over $V$ at scales tending to zero would have a subsequential limit in
$(X_0\setminus U)\cap\pi^{-1}(\overline V)$, a contradiction.
A finite outer subcover gives the common minimum scale cutoff and
maximum bounds. The same argument applies to the one-copy graph.
\end{proof}

\begin{lemma}[An endpoint cutoff chosen after contraction]
\label{rw:fibre:candidate-cover}
There is $p_\sharp>0$, chosen before the primary scale cutoff,
and finitely many local candidate neighborhoods that contain every
counted cycle with $p\ge p_\sharp/16$ at sufficiently small
original parameter. The regular member contains all sufficiently
small endpoint data with
$p_\sharp/16\le p\le p_\sharp/2$.
Each other member is either separated from its local center set
on its closure or has a common center-comparison domain and a
positive contraction exponent $b_j>p_\sharp^2$.
Secondary center estimates retain any prescribed finite power
of $\lambda$, or of the wedge height, before division.
\end{lemma}
\begin{proof}
Start with the nonregular part of the compact refined limit set
in Lemma~\ref{rw:fibre:candidate-space}. At a noncenter point
shrink the closed coefficient neighborhood away from the center
set. At a center point the classification gives a divided regular,
two-root, finite-saddle, weak, focus or joint chart.
There are no additional unprepared center losses: the exact
quasi-invariant \eqref{rw:finite:quasi-invariant} implies that an
incoming finite root on a center also has an outgoing root; a
one-infinite center has
$\bar a=-\bar c,\bar m=-\bar B$, so its finite lower factor
equals the upper entry factor; and a lower double root in the
secondary corner has $|\bar a|=2\sqrt{\bar B}>0$ and is off
the secondary center. The repeated coincidence is precisely the
joint chart.

On each center neighborhood, first retain a fixed lower subsegment
with positive action. After reducing the coefficient neighborhood,
normal contraction on this subsegment bounds the absolute coefficient
derivatives of the complete central deviations by
$C\epsilon^{-N}e^{-b_j/\epsilon}$, with $b_j>0$.
These are estimates of the absolute deviations; the logarithm of an
input-independent slow contraction may diverge. The large-power
preparation extends the same estimate by zero up to a lost exit.
Section~5 supplies the real center and the whole input domain on which
to divide. Double division uses two ordinary coefficient derivatives
of the absolute difference, and a single component uses one. The
budgets retained in Sections~7--9 include these derivatives; this step
takes no further input derivative.
In the corner use the weighted absolute estimates of
Lemma~\ref{rw:secondary:flat-budget}. For the present argument
retain power two before the coefficient division; if a later local
coordinate requires more, retain its specified finite power first.
The exact exponent in Proposition~\ref{rw:exact-long} makes
these choices uniform without bounding $\epsilon|\log\lambda|$.
For the wedge choose its fixed widths $|\gamma|,|J|<\delta$
inside the regular-lower center box whenever $|\bar a|$ is
small. The remaining wedge data have $|\bar a|$ bounded
below and are primary off-center data. Thus an off-center
secondary direction attached to a first wedge exit is not
used to estimate that exit by a power of $\lambda$; its
actual radial power is a power of $H$.

Take a finite cover of this refined nonregular set, and then choose
$p_\sharp\le\chi/2$ so that $p_\sharp^2<\min_j b_j$,
and decrease it to meet the small-endpoint error bounds in the
source and regular-column preparations. These choices precede
the smallness choice for $\epsilon$.
The compact complement of the chosen open neighborhoods is
strictly regular. Enlarge that compact regular set to contain
$p_\sharp/16\le p\le p_\sharp/2$ and the limits of cycles
in the two initial energy classes with $p\ge p_\sharp/16$.
The latter enlargement is regular: the source endpoints satisfy
$p=\sqrt{2\epsilon h e_+}\le2\sqrt\eta$ and
$q/p=\sqrt{e_-/e_+}=1+O(\eta)$, while the lower and upper
thresholds have a fixed positive lower bound. The compact-energy
class has $p,q\to0$ and causes no further endpoint limit.
The regular upper return and the positive lower polynomials
therefore give a common compact tube for this enlargement.

Apply Lemma~\ref{rw:fibre:compact-selection} to this one-copy
refined graph. Positive endpoint limits from the initial classes lie
in the regular enlargement, with $q$ bounded below when
$p\ge p_\sharp/16$ by the same source estimates. Thus the selected
finite family covers all the asserted cycle data.
\end{proof}

The candidate cover locates each outer closing orbit. To obtain one
sharp count, we next extract its coefficient inequality; saturation
then supplies all intervening comparison paths.

\begin{proposition}[Forcing supplied by an outer closing zero]
\label{rw:outer-forcing}
Choose the field orientation with $\ell\ge0$.
After a common original parameter shrink, a counted cycle with
$p\ge p_\sharp$ either lies in a constant-sign local comparison
whose full remaining physical passing interval is covered by that
same comparison, or supplies the coefficient inequality
\begin{equation}\label{rw:fibre:original-forcing}
 |\ell|\le e^{-p_\sharp^2/(20\epsilon)}
                    (|\tau_0|+|\sigma_0|).
\end{equation}
Equivalently,
\begin{equation}\label{rw:fibre:normalized-forcing}
 |\nu|\le e^{-p_\sharp^2/(20\epsilon)}
                    (|\tau_1|+\epsilon|\sigma_1|).
\end{equation}
An interior constant-sign piece is not sufficient for the first
alternative. On a center fibre with a closing zero there are no
isolated closing zeros.
\end{proposition}
\begin{proof}
Apply the candidate cover to the outer cycle. A constant-sign chart
containing the entire remaining passing interval gives the first
alternative. In every other chart we derive the coefficient inequality,
including on constant-sign pieces separated from the center set.

The nonnegative lower cone and
\eqref{rw:finite:action-budget}, applied at $p_\sharp$, give
\[
 s\le \exp[-p_\sharp^2/(5\epsilon)].
\]
Indeed
$A_+(p_\sharp)\ge
p_\sharp^2/[2(1+p_\sharp+p_\sharp^2)]$ and the error in
the budget is absorbed after $p_\sharp$ is fixed.
Since $|\ell|<3s$,
any closed primary neighborhood on which
$|\tau_1|+|\sigma_1|\ge g_0>0$ satisfies
\eqref{rw:fibre:original-forcing}: use
$|\tau_0|+|\sigma_0|\ge\rho^3g_0$ and absorb this fixed
power of $\rho$ in the strict gap between the two exponents.

At an off-center secondary first exit, the primary action estimate
alone does not retain the vanishing secondary weight. We use the exact
long match to recover it.
Suppose first $|C|\ge g_0/2$, where
$|\bar a|+|C|\ge g_0$ on the chosen closed neighborhood.
The exact long match gives an outgoing horizontal input
$u_+\le C_1\lambda^{k/b}$, with
$b=1-\epsilon\bar B$, $k=1-2\epsilon\bar B$.
The lower common power has exponent $b/(\epsilon\bar B)$.
The bounded scaled action and scaled prefactor, including a
possible additional contracting lower double passage, imply
\begin{equation}\label{rw:fibre:weighted-budget}
 s\le I_{\rm out}
       \le e^{C_0/\epsilon}\lambda^{k/(\epsilon\bar B)}.
\end{equation}
The first inequality uses the actual reference
$s_+=-\ell H_+\le0$. Choose $\kappa_*>0$ with
$k/\bar B\ge\kappa_*$, then $\lambda_0>0$ so that
\[
 (\kappa_*/2)\log(1/\lambda_0)
                    \ge C_0+p_\sharp^2/5,
 \qquad \epsilon\le\kappa_*/4.
\]
For $\lambda\le\lambda_0$,
\eqref{rw:fibre:weighted-budget} yields
$s\le\lambda^2e^{-p_\sharp^2/(5\epsilon)}$.
For $\lambda\ge\lambda_0$ the primary budget gives the same
bound with the fixed factor $\lambda_0^{-2}$. Consequently
$|\nu|\le\lambda|C|/2$ at small scale and
\[
 |\bar c|=|\lambda C-\nu|\ge\lambda g_0/4,\qquad
 |\sigma_0|\ge c\rho^3\lambda g_0 .
\]
The extra factor $\lambda^2$ proves
\eqref{rw:fibre:original-forcing} uniformly, even if
$\epsilon|\log\lambda|$ is arbitrarily large.
If $|C|<g_0/2$, then $|\bar a|\ge g_0/2$; the corner
has $|\bar c|<g_0/4$, so this reduces to the primary
off-center estimate.

On a retained center neighborhood, cancellation on the real center
recovers the generator factor. Write the complete closing function as
$Z-\rho\nu H$, with $H>h_*$ independent of the section input.
The actual center division gives
\[
 |Z|\le C\epsilon^{-N}e^{-b_j/\epsilon}
             (|\tau_1|+|\sigma_1|+|\nu|).
\]
The local inverse charts and the two ideal comparisons in
Proposition~\ref{rw:generator-division} justify the indicated
generators. In secondary variables retain the weighted version
before returning to primary coefficients: a term
$\lambda^2|C|$ is bounded by
$|\bar d|=\lambda|C|$, and the wedge version is identical
with $H^2|\gamma|\le|\bar d|$ and its retained height.
The small-$|\bar a|$ wedge box was included in these center
comparisons independently of its secondary direction. Since
$|\bar d|\le|\bar c|+|\nu|$ and
$\bar B+\bar m>1/8$ here, the resulting bound is again the
displayed primary one. No inverse power of $\lambda$ is
introduced.
At a closing zero Lemma~\ref{rw:centers:relative-forcing}
absorbs the $\nu$ term. Since $b_j>p_\sharp^2$, the fixed
powers of $\epsilon$ needed to return to original generators
are absorbed with room to spare.

Finally consider the compact regular member at $p\ge p_\sharp$.
Retain a slightly larger actual tube with $p\ge p_\sharp/2$.
The regular upper map has $q\ge p(1-p)$ at small endpoints;
hence both retained actions have a lower bound greater than
$p_\sharp^2/10$ after $p_\sharp$ is reduced. The regular
absolute maps and double center division yield
\[
 |\ell|\le C\epsilon^{-N}e^{-p_\sharp^2/(10\epsilon)}
                      (|\tau_0|+|\sigma_0|),
\]
which implies the required exponent $p_\sharp^2/20$.
The finite candidate cover makes these choices uniform. If the
generators vanish, absorption first gives $\nu=0$, and the actual
center identity excludes isolated zeros.
\end{proof}
